\label{sec:third}
This subsection makes Theorem~\ref{th:cycles} explicit for $j=2$ and one profile. Here $f=m_1(\mu)$ and $c=m_2(\mu)$ count fixed points and $2$-cycles, and the third term involves exactly these two numbers.
\begin{theorem}\label{th:third}
For $d\ge5$, the third coefficient in Theorem~\ref{th:cycles} is
\begin{equation}\label{eq:third}
 C_2(t)=A_{d-3}(t)\bigl[a(1+t)^2-Qt\bigr],
\end{equation}
where
\begin{align}
 a&=\binom{d+f-2}{2}-c,\label{eq:third-a}\\
 Q&=\frac{(d-1)\bigl[2c(d-2)+(d-3)(d-4)+2f(d+f-5)\bigr]}
 {(d-4)(d-2)}.\label{eq:third-Q}
\end{align}
Thus, on the admissible parity,
\begin{align*}
 \frac{d!}{2}h_r(\mu;t)
 ={}&(d-1)^rA_{d-1}(t)
 -(d+f-2)(1+t)(d-2)^rA_{d-2}(t)\\
 &+(d-3)^rA_{d-3}(t)\bigl[a(1+t)^2-Qt\bigr]
 +O_{d,\mu}\bigl((1+r)^d(d-4)^r\bigr).
\end{align*}
\end{theorem}
\begin{proof}
Put $\ell=d-3$ and $D=d-1$. At the pole $z=1/\ell$ abbreviate
\[
 F_b=\frac{1+(t-1)b/\ell}{1-b/\ell}.
\]
The generating function for connected words is given by
Lemma~\ref{lem:mobius} and the content product \eqref{eq:char}. A block supporting the pole has size at least $d-2$.
Since $d\ge5$, there is a unique such block, and the pole is simple.
In each contributing diagram the contents $0,1,\ldots,\ell$ supply
the common polar coefficient $A_\ell(t)$. We compute the remaining
multiplier after multiplying the residue by $d!$.

For the one-block partition, the four possible shapes and their
contributions are
\begin{align*}
 (d)       &\quad F_DF_{D-1},\\
 (d-1,1)   &\quad D(f-1)F_{D-1}F_{-1},\\
 (d-2,2)   &\quad \frac{d(d-3)}{2}
                   \left(\binom f2+c-f\right)F_{-1},\\
 (d-2,1,1) &\quad \frac{(d-1)(d-2)}{2}
                   \left(\binom{f-1}{2}-c\right)F_{-1}F_{-2}.
\end{align*}
The character identities in the last two rows follow directly
from Jacobi--Trudi. The identity
$s_{(d-2,2)}=h_{d-2}h_2-h_{d-1}h_1$
counts invariant two-element sets and subtracts fixed points.
The expansion
\[
 s_{(d-2,1,1)}
 =h_{d-2}(h_1^2-h_2)-h_{d-1}h_1+h_d
\]
has character $f(f-1)-\binom f2-c-f+1
=\binom{f-1}{2}-c$.
Here $\binom x2=x(x-1)/2$, including $x=-1$ when $f=0$.
Their values at the identity give the dimensions shown above.

For a singleton and its complement, the singleton must be one of
the $f$ fixed points. The Möbius coefficient is $-1$. The
complementary diagrams with content $\ell$ are $(d-1)$ and
$(d-2,1)$, and the restricted permutation has $f-1$ fixed points.
After the factor $d!/(d-1)!=d$, the contribution is
\[
 -df\bigl[F_{D-1}+(d-2)(f-2)F_{-1}\bigr].
\]
When $f=0$ this term is zero.

The remaining possibilities have two letters outside the large
block. A two-element block is either a pair of fixed points, of
which there are $\binom f2$, or the support of a $2$-cycle, of
which there are $c$. Its unconstrained word generating function
is $(F_1+F_{-1})/2$ in the first case and
$(F_1-F_{-1})/2$ in the second. Both partitions have Möbius
coefficient $-1$. With $u=f(f-1)$ and
$d!/(d-2)!=dD$, their contribution is
\[
 -\frac{dD}{4}\bigl[(u+2c)F_1+(u-2c)F_{-1}\bigr].
\]
Finally, two singleton blocks require two fixed points.
There are $\binom f2$ such partitions, each with Möbius coefficient
$+2$, giving $dD u$. This exhausts all set partitions capable
of supporting the pole.

Let $W(t)$ be the sum of the seven contributions just listed.
The values of the remaining linear factors are
\[
 F_D=1-\frac{d-1}{2}t,\quad F_{D-1}=1-(d-2)t,\quad
 F_{-1}=1-\frac{t}{d-2},\quad
 F_{-2}=1-\frac{2t}{d-1},\quad F_1=1+\frac{t}{d-4}.
\]
Their substitution gives
$[t^0]W=[t^2]W=a$ and $[t]W=2a-Q$,
with $a,Q$ as in \eqref{eq:third-a}--\eqref{eq:third-Q}.
Consequently $C_2=A_\ell W$ has the stated form.
The negative pole contributes equally on admissible lengths,
and all remaining poles have reciprocal magnitudes at most
$d-4$. Partial fractions give the asserted remainder.
\end{proof}

For $d=5$, the full-cycle profile has $a=3$ and $Q=8/3$, giving
$C_2(t)=t(1+t)\bigl[3(1+t)^2-\frac83t\bigr]$.
For $d=6$ and $\mu=(2,2,2)$, the values are $a=3$ and $Q=75/4$.
The requirement $d\ge5$ ensures that $F_1$ is finite at the
third pole and that the large block is unique; for $d=4$ the third pole is $z=1$, where $F_1$ itself has a pole.
